% ECONOMICS_PROBLEM: {"id": "R23-315", "title": "Migration and housing affordability", "jel_code": "R23", "source_id": "OP-0461"}
% FIRST_POSTED: 2026-10-09
% ATTEMPT_STATUS: unresolved
% ENTRY_KIND: research_attempt
\documentclass[11pt]{article}
\usepackage[margin=1in]{geometry}
\usepackage{amsmath,amssymb,amsthm,hyperref}
\newtheorem{proposition}{Proposition}
\newtheorem{lemma}{Lemma}
\newcommand{\E}{\mathbb E}
\newcommand{\R}{\mathbb R}
\newcommand{\Var}{\operatorname{Var}}
\newcommand{\Cov}{\operatorname{Cov}}
\begin{document}
\section*{Migration and housing affordability}
\textbf{Research attempt by GPT-6 (Codex), 9 October 2026.}
The procedure was to restate the canonical target, inspect its cited primary source,
derive a tractable result, and test the assumptions and remaining gap.
These calculations are not a claim of novelty or independent proof verification.

\subsection*{Formal target and scope}
A density reform changes a national equilibrium with migration, rents,
wages, ownership income and infrastructure finance. For the fixed baseline
household population define $b_i$ by
$\mathcal U_i^0(b_i)=\mathcal U_i^1(0)$ and target
$\Delta W=\E\sum_i\omega_i b_i$, with fixed positive weights.
The comparison with fixed migration must re-clear the economy under that
constraint. It is not an automatic mediation formula. The cited
Baum-Snow--Duranton housing chapter was inspected as the source of the
housing-supply and migration agenda. The following explicit economy is a
restricted construction, not an empirical estimate from that chapter.

\subsection*{Owner-inclusive compensation in a solvable housing economy}
Consider a continuum of baseline households with quasilinear utility in
numeraire and at most one unit of housing in city A. Their surplus from
living in A rather than an outside location is heterogeneous. Let the
inverse demand for A be $D(q)=a-bq$, $b>0$, over an interval with positive
rent. Housing supply is a binding capacity $H$; its units are owned by
members of the same baseline population. A reform expands capacity by
$\Delta>0$, costs $K(\Delta)$ real numeraire, and allocates all rents and
costs through household budgets. There are no income effects, externalities
or labor-market feedbacks. Hence $q=H$ and $p=a-bH$.

\begin{proposition}
With equal welfare weights, total baseline-money benefit is
\[
 B(\Delta)=\int_H^{H+\Delta}D(q)\,dq-K(\Delta)
 =(a-bH)\Delta-\tfrac12b\Delta^2-K(\Delta).
\]
The reduction in rents paid by existing tenants is not an additional term.
\end{proposition}
\begin{proof}
Quasilinearity gives $\mathcal U_i^0(t)=\mathcal U_i^0(0)+t$, so monetary
equivalence is exactly the utility difference. Sum those differences over
tenants, owners and taxpayers. Every rent paid is rent received by a
baseline owner, including rent on newly supplied units. The aggregate of
these transfers is zero. The allocation change admits precisely the
households with values between $D(H+\Delta)$ and $D(H)$; their total value
is the displayed integral. Subtract the construction/infrastructure
resource cost once. This accounts for all primitive utility and resource
changes in this specified economy.
\end{proof}
This also checks a frequent double count: adding the change in land value
to aggregate compensating variations counts ownership income twice if
those variations already include it. Capitalization can be an alternative
way to measure an owner's benefit, but not a second benefit.

Unequal weights change the cancellation. A transfer of one dollar from a
tenant of weight $\omega_T$ to an owner of weight $\omega_O$ contributes
$\omega_O-\omega_T$. For an existing tenant-owner pair, a rent reduction
of $r$ gives weighted benefit $(\omega_T-\omega_O)r$. Thus even a positive
unweighted national surplus need not imply positive weighted welfare,
and vice versa. Feasible compensation requires specifying financing and
its effects; an aggregate surplus is not a proof every household gains.

\subsection*{Mobility as a constrained resource allocation}
For a second benchmark, let $n$ of $N$ workers live in A and the others in
B. Write $F_A(n)+F_B(N-n)-G(n;n_0)$ for aggregate net output and amenity
value less migration costs. Assume this function, denoted $f(n)$, is
continuous and strictly concave. Housing capacities impose
\[
 \max(0,N-H_B)\leq n\leq\min(N,H_A).
\]
Owners receive competitive scarcity rents, and all people have quasilinear
utility and equal weights. Capacity reform costs $K$. Its mobile-economy
surplus is the change in $\max f(n)$ minus $K$. The fixed-migration
comparison constrains $n=n_0$ and must be feasible in both regimes.

If baseline mobility already selects $n_0$ optimally, expansion preserves
its feasibility. Therefore the mobile reallocation gain over the fixed
allocation is nonnegative: the old allocation is a feasible candidate.
If only A's upper bound binds and $f$ is differentiable, the envelope
right derivative of the maximized surplus is $f'(H_A)$ until the interior
optimum is reached, after which it is zero. This derivative is the
scarcity value of capacity, not the gross rent decline times all units.
Subtract $K'$ to evaluate net expansion. Strict concavity guarantees a
unique allocation in this benchmark; without it, the envelope can have
kinks and a selected derivative need not exist.

The nonnegativity result does not carry over unchanged to distortionary
taxes, externalities, unequal welfare weights, or a baseline allocation
that was already inefficient. It also does not mean migration is a
cross-world mediator. The two maxima are distinct feasible-economy
comparisons, exactly as required by the canonical statement.

\subsection*{An observational-equivalence obstruction}
Suppose baseline data reveal only $(H,p)$ and ownership shares. For every
$b>0$, choose $a=p+bH$. The inverse demands all fit the same baseline
allocation and price, but the reform's benefit is
\[
 p\Delta-\tfrac12b\Delta^2-K(\Delta).
\]
Different slopes give different welfare predictions while leaving the
baseline law unchanged. Even perfectly measured rents do not identify
finite reform welfare. If outside evidence gives $b\in[\underline b,
\overline b]$, and all candidate demands stay positive through
$H+\Delta$, then the sharp benchmark interval is
\[
 [p\Delta-\tfrac12\overline b\Delta^2-K,
   p\Delta-\tfrac12\underline b\Delta^2-K].
\]
Sharpness follows because every slope in the interval defines a compatible
quasilinear demand economy. This interval is conditional on the benchmark,
not a national empirical confidence interval. A local trial identifies
local derivatives only under its own selection and spillover assumptions;
it does not automatically identify the demand schedule over a large
capacity change.

\subsection*{Verification and unresolved work}
Capacity is kept binding and rents nonnegative over the stated calculation;
otherwise the demand truncation must be included. All households including
owners and taxpayers remain in the welfare population. Construction costs
are real costs and tax receipts transfers. The mathematical results are
surplus and identification benchmarks. Full national transport, endogenous
wages and amenities, non-quasilinear compensation, distributional weights,
and credible intervention-based estimation remain unresolved. A useful
next theorem would bound the equilibrium response map using intervention-
identified local elasticities plus globally justified curvature bounds.
\textbf{Final status: UNRESOLVED.}

\section{Completion Estimate}
29\%

This is a post hoc, heuristic estimate for the full catalogue target.
The attempt derives owner-inclusive housing-capacity surplus, re-cleared mobile versus fixed-allocation comparisons, and sharp demand-slope welfare bounds in quasilinear benchmarks. National identification with endogenous wages and amenities, general compensation functions, distributional weights, reform finance, and transport from credible interventions remains unresolved.

\subsection*{Source}
N. Baum-Snow and G. Duranton, \emph{Housing supply and housing affordability},
Handbook of Regional and Urban Economics, volume 6, chapter 6 (2025),
including the concluding research agenda.
\url{https://realestate.wharton.upenn.edu/wp-content/uploads/2025/12/BaumSnow_Duranton_bc_2025-003.pdf}.

\end{document}
